% ============================================================
\chapter*{Preface}
\addcontentsline{toc}{chapter}{Preface}
% ============================================================

\section*{Course Objectives}

Convex optimization occupies a central place in modern applied mathematics.
It provides a rigorous theoretical framework and efficient algorithms for solving
a broad class of optimization problems arising in machine learning, signal processing,
finance, operations research, and many other fields.

This Master-level course is intended for students with solid backgrounds in real
analysis, linear algebra, and topology. It covers the theoretical foundations of
convexity, optimality conditions, duality theory, and the most important algorithmic
methods used in practice.

The approach adopted is both mathematically rigorous and application-oriented.
Each theoretical concept is accompanied by geometric interpretations, concrete
examples, and Python implementations.

\section*{Course Organization}

The course is organized into ten chapters, structured in a logical progression:

\begin{enumerate}
    \item \textbf{Convex Sets} --- Fundamental definitions, topological properties,
          operations preserving convexity, cones, polyhedra, separation theorems.
    \item \textbf{Convex Functions} --- Differential characterizations,
          fundamental inequalities (Jensen, gradient), continuity and differentiability.
    \item \textbf{Subdifferentials and Subgradients} --- Nonsmooth convex analysis,
          subdifferential calculus, Moreau--Rockafellar rules.
    \item \textbf{Fenchel--Moreau--Rockafellar Duality} --- Convex conjugation,
          biconjugation theorem, Fenchel duality.
    \item \textbf{Optimality Conditions --- KKT} --- Necessary and sufficient
          conditions, constraint qualifications.
    \item \textbf{Lagrangian Duality} --- Dual problem, duality gap,
          strong duality, Slater's theorem, economic interpretation.
    \item \textbf{Gradient Descent and Variants} --- Fixed and adaptive step-size,
          Nesterov's accelerated gradient, stochastic gradient.
    \item \textbf{Proximal Methods and ADMM} --- Proximal operator,
          ISTA/FISTA algorithms, ADMM, decomposition.
    \item \textbf{Interior Point Methods} --- Barrier functions,
          central path method, polynomial complexity.
    \item \textbf{Applications} --- LASSO, SVM, logistic regression,
          portfolio optimization.
\end{enumerate}

\section*{Prerequisites}

The essential prerequisites for this course are:

\begin{itemize}
    \item \textbf{Real Analysis} --- Sequences, series, continuity, differentiability
          in $\R^n$, convergence theorems.
    \item \textbf{Linear Algebra} --- Vector spaces, matrices, eigenvalues,
          matrix decompositions (SVD, Cholesky), quadratic forms.
    \item \textbf{Topology} --- Open and closed sets, compactness, connectedness
          in $\R^n$, notion of relative interior.
    \item \textbf{Programming} --- Basic knowledge of Python (NumPy, matplotlib).
          The \texttt{cvxpy} library will be used for implementations.
\end{itemize}

\section*{Notation}

We use the following notation throughout the course:

\begin{center}
\renewcommand{\arraystretch}{1.4}
\begin{tabular}{cl}
    \toprule
    \textbf{Notation} & \textbf{Description} \\
    \midrule
    $\R^n$ & Euclidean space of dimension $n$ \\
    $\inner{x}{y}$ & Inner product $\sum_{i=1}^n x_i y_i$ \\
    $\norm{x}$ & Euclidean norm $\sqrt{\inner{x}{x}}$ \\
    $\norm{x}_1$ & $\ell_1$ norm: $\sum_i \abs{x_i}$ \\
    $\norm{x}_\infty$ & $\ell_\infty$ norm: $\max_i \abs{x_i}$ \\
    $B(x,r)$ & Open ball centered at $x$ with radius $r$ \\
    $\overline{C}$ & Closure of the set $C$ \\
    $\mathrm{int}(C)$ & Interior of the set $C$ \\
    $\mathrm{ri}(C)$ & Relative interior of $C$ \\
    $\mathrm{aff}(C)$ & Affine hull of $C$ \\
    $\mathrm{conv}(S)$ & Convex hull of $S$ \\
    $\mathrm{cone}(S)$ & Conic hull of $S$ \\
    $\mathrm{dom}(f)$ & Effective domain of $f$ \\
    $\mathrm{epi}(f)$ & Epigraph of $f$ \\
    $f^*$ & Fenchel conjugate of $f$ \\
    $\partial f(x)$ & Subdifferential of $f$ at $x$ \\
    $\nabla f(x)$ & Gradient of $f$ at $x$ \\
    $\nabla^2 f(x)$ & Hessian of $f$ at $x$ \\
    $\argmin_{x} f(x)$ & Set of minimizers of $f$ \\
    $A \succeq 0$ & $A$ is positive semidefinite \\
    $A \succ 0$ & $A$ is positive definite \\
    $\mathbb{S}^n_+$ & Cone of $n \times n$ PSD matrices \\
    $\iota_C$ & Convex indicator function of $C$ \\
    $\sigma_C$ & Support function of $C$ \\
    $\mathrm{prox}_f$ & Proximal operator of $f$ \\
    \bottomrule
\end{tabular}
\end{center}

\section*{Conventions}

\begin{itemize}
    \item Unless otherwise stated, all spaces considered are finite-dimensional over $\R$.
    \item Convex functions considered take values in $\R \cup \{+\infty\}$
          (proper convex functions).
    \item The symbol $\square$ or $\qed$ marks the end of a proof.
    \item Exercises are classified by increasing difficulty:
          $(\star)$ basic, $(\star\star)$ intermediate, $(\star\star\star)$ advanced.
    \item \texttt{algorithme} blocks present pseudo-code for the methods studied.
    \item \texttt{codeblock} blocks contain executable Python code.
\end{itemize}

\section*{Guiding Thread}

The guiding thread of this course is the general convex optimization problem:
\[
    \min_{x \in \R^n} \; f(x) \quad \text{subject to} \quad
    g_i(x) \leq 0, \; i=1,\ldots,m, \quad
    h_j(x) = 0, \; j=1,\ldots,p,
\]
where $f$ and the $g_i$ are convex and the $h_j$ are affine. We study successively:
\begin{itemize}
    \item the \emph{geometry} of the feasible set (Chapters 1--2),
    \item the \emph{analysis} of objective functions (Chapters 2--4),
    \item \emph{optimality conditions} and \emph{duality} (Chapters 5--6),
    \item \emph{algorithms} for solving (Chapters 7--9),
    \item concrete \emph{applications} (Chapter 10).
\end{itemize}

\section*{Pedagogical Approach}

Each chapter follows a similar structure:
\begin{enumerate}
    \item \textbf{Motivation} --- Why this topic matters.
    \item \textbf{Theory} --- Definitions, theorems with complete proofs.
    \item \textbf{Geometric Interpretation} --- TikZ illustrations and intuitions.
    \item \textbf{Examples} --- Detailed examples and counterexamples.
    \item \textbf{Implementation} --- Python/cvxpy code.
    \item \textbf{Exercises} --- Problems of varying difficulty.
\end{enumerate}

\section*{Main References}

\begin{enumerate}
    \item S.~Boyd and L.~Vandenberghe, \emph{Convex Optimization},
          Cambridge University Press, 2004.
    \item R.T.~Rockafellar, \emph{Convex Analysis},
          Princeton University Press, 1970.
    \item J.-B.~Hiriart-Urruty and C.~Lemar\'echal,
          \emph{Fundamentals of Convex Analysis}, Springer, 2001.
    \item Y.~Nesterov, \emph{Introductory Lectures on Convex Optimization},
          Springer, 2004.
    \item A.~Beck, \emph{First-Order Methods in Optimization}, SIAM, 2017.
    \item N.~Parikh and S.~Boyd, \emph{Proximal Algorithms},
          Foundations and Trends in Optimization, 2014.
    \item H.H.~Bauschke and P.L.~Combettes,
          \emph{Convex Analysis and Monotone Operator Theory in Hilbert Spaces},
          Springer, 2017.
    \item D.P.~Bertsekas, \emph{Convex Optimization Algorithms}, Athena Scientific, 2015.
\end{enumerate}

\vfill

\begin{flushright}
\textit{The Author}\\
March 2026
\end{flushright}
